\section{Data and Methods}
\label{sec:data_methods}

The methodological workflow is summarised in \cref{fig:workflow} and proceeds in
five stages: (1) flood inventory construction; (2) conditioning factor assembly
and multicollinearity screening; (3) watershed graph construction;
(4) model training and spatial block cross-validation; and
(5) conformal prediction and uncertainty quantification.

\subsection{Flash Flood Inventory}
\label{sec:inventory}

\subsubsection{SAR-Based Detection (Primary)}
\label{sec:sar_inventory}

We processed all available Sentinel-1 C-band SAR Ground Range Detected (GRD) scenes
over HP from January 2018 to December 2024 using the Google Earth Engine (GEE)
cloud computing platform \citep{gorelick2017gee}.
Sentinel-1 SAR at 5.4\,GHz C-band penetrates monsoon cloud cover,
enabling flood detection during precisely the peak flash flood season
when optical imagery is systematically unavailable
\citep{nagamani2024sar, bentivoglio2022review}.

The detection methodology followed established SAR change-detection protocols
\citep{martinis2018sar}:
(1) pre-event and post-event imagery were composited to 10-day medians
to reduce noise;
(2) VV-polarisation backscatter was used (more stable over water than VH);
(3) a $-3$\,dB change threshold identified pixels with significant backscatter
reduction, indicative of specular reflection from inundated surfaces;
(4) a focal-mode filter ($3{\times}3$ kernel) removed isolated pixels;
(5) the JRC Global Surface Water mask \citep{pekel2016jrc} removed
permanent water bodies; and
(6) flood patches smaller than 0.54\,ha (6 connected pixels at 30\,m)
were discarded as probable noise.
Each detected event was assigned a date, trigger classification
(monsoon/cloudburst/GLOF/uncertain), and basin attribution.

The resulting inventory yielded XX flood inundation polygons spanning
XX event-days from 2018 to 2024.
Flood occurrence points were extracted at a density of one per 0.09\,km$^2$
(one per 30\,m pixel) within each inundation polygon,
subject to a minimum inter-point distance of 90\,m to avoid spatial
pseudo-replication within a single flood event.

\subsubsection{Historical Records (Supplementary)}
\label{sec:historical_inventory}

SAR-detected flood points were supplemented with:
(1) the HiFlo-DAT database \citep{hiflodot2025} (Kullu district, 128 events,
1900--2023; only events post-2010 with adequate location precision were used);
(2) NDMA situation reports for 2018--2024, georeferenced to reach-level
using NH kilometre posts and district revenue maps;
(3) digitised flood extents from \citet{inventories2020beas}
for Beas tributary channels.

Combined, the inventory contains XX flood occurrence points
(training: 2018--2022, N\,=\,XX; temporal test: 2023, N\,=\,XX).

\subsubsection{Non-Flood Point Sampling}
\label{sec:nonfloods}

Non-flood (absence) points were generated by stratified random sampling
within the HP boundary, subject to: (1) minimum 1,000\,m Euclidean distance
from any flood occurrence point, to prevent spatial leakage of positive class
information into the negative class; and (2) proportional spatial coverage
of all physiographic zones.
A 5:1 negative-to-positive ratio was maintained following recommendations
by \citet{dataquality2025}.

\subsubsection{Temporal Split}
\label{sec:temporal_split}

Events from 2018--2022 form the training set.
Events from July--September 2023 form an independent temporal test set.
This split tests temporal transferability of susceptibility models
\citep{ghosh2023karnali} and exploits the 2023 season's exceptional documentation
quality as a held-out benchmark.

\subsection{Conditioning Factors}
\label{sec:factors}

Thirteen conditioning factors were assembled at 30\,m resolution,
reprojected to UTM Zone 43N (EPSG:32643), and clipped to HP extent.
\cref{tab:factors} summarises each factor and its data source.

\begin{table}[htbp]
  \centering
  \caption{Conditioning factors, data sources, and descriptive statistics.
    All factors resampled to 30\,m resolution and reprojected to UTM Zone 43N.}
  \label{tab:factors}
  \small
  \begin{tabularx}{\textwidth}{lllX}
    \toprule
    Factor & Unit & Source & Rationale \\
    \midrule
    Elevation              & m         & Copernicus GLO-30    & Controls temperature, snowmelt contribution, flood transmission distance \\
    Slope                  & degrees   & Derived from DEM     & Primary runoff velocity and infiltration control \\
    Aspect                 & degrees   & Derived from DEM     & Snow-shading and moisture retention effects \\
    Plan curvature         & m$^{-1}$  & Derived from DEM     & Flow convergence / divergence \\
    Profile curvature      & m$^{-1}$  & Derived from DEM     & Flow acceleration \\
    TWI                    & --        & D8 flow accumulation & Topographic Wetness Index; soil moisture proxy \\
    SPI                    & --        & D8 flow accumulation & Stream Power Index; erosion capacity \\
    TRI                    & m         & Derived from DEM     & Terrain Ruggedness Index; surface roughness \\
    Mean annual rainfall   & mm/year   & GPM-IMERG v07        & Long-term moisture availability \\
    Extreme rainfall (p95) & mm/day    & GPM-IMERG v07        & High-intensity event triggering threshold \\
    LULC                   & class     & ESA WorldCover 2021  & Land surface hydrological response \\
    Soil clay content      & \%        & SoilGrids v2.0       & Soil infiltration capacity \\
    Distance to river      & m         & OSM + DEM drainage   & Proximity to flood-prone channels \\
    \bottomrule
  \end{tabularx}
\end{table}

\subsubsection{Terrain Factor Derivation}
\label{sec:terrain}

The Copernicus GLO-30 digital elevation model (30\,m, released 2021) was
used as the primary elevation source, preferred over SRTM for its improved
performance in high-relief terrain \citep{fabdem2022}.
Slope, aspect, plan curvature, and profile curvature were computed using
standard finite-difference algorithms.
TWI and SPI were derived from D8 flow accumulation computed with \textit{pysheds}
\citep{pysheds2020}:
\begin{equation}
  \mathrm{TWI} = \ln\!\left(\frac{A_s}{\tan\beta}\right)
  \label{eq:twi}
\end{equation}
\begin{equation}
  \mathrm{SPI} = A_s \cdot \tan\beta
  \label{eq:spi}
\end{equation}
where $A_s$ is the specific catchment area (m$^2$\,m$^{-1}$) and
$\beta$ is local slope in radians.
TRI was computed as the mean absolute elevation difference between each cell
and its eight neighbours \citep{riley1999tri}.

\subsubsection{Rainfall Factors}
\label{sec:rainfall}

Monthly GPM-IMERG Final Run v07 precipitation composites (0.1° resolution)
were downloaded for 2000--2023 and spatially interpolated to 30\,m
using cubic spline resampling.
Two rainfall factors were computed: (1) mean annual precipitation (mm/year)
as a long-term climate proxy; and (2) the 95th percentile of daily event
precipitation (p95, mm/day) as a trigger intensity proxy.

\subsubsection{Multicollinearity Assessment}
\label{sec:multicollinearity}

Pearson pairwise correlation was computed for all 13 continuous factors
on a stratified random sample of 10,000 pixels.
Pairs with $|r| > 0.80$ were flagged for removal.
Variance Inflation Factor (VIF) was computed for the retained factor set,
with VIF\,$>$\,10 as the removal threshold.
The final retained factor set is reported in Section~\ref{sec:factors_retained}.

\subsection{Watershed Graph Construction}
\label{sec:graph}

To encode upstream--downstream hydrological connectivity for GNN training,
we delineated sub-watershed units from the GLO-30 DEM using the
D8 contributing area algorithm with a minimum catchment area threshold
of 150\,km$^2$, yielding XX sub-watersheds.
Larger catchments (lower minimum area) were avoided to prevent
coarse spatial resolution that obscures localized flash flood processes.

Directed edges were placed from each sub-watershed to all sub-watersheds
directly downstream (i.e., receiving outflow from the source watershed),
inferred from the DEM-derived drainage network.
For GNN message-passing, reverse (undirected) edges were also added
to allow upstream propagation of information during training
\citep{hamilton2017graphsage}.
Edge weights were proportional to the upstream catchment area ratio,
reflecting the discharge contribution of each upstream node.

Node features were computed as the area-weighted mean of each conditioning factor
over the sub-watershed polygon, yielding a feature vector of dimensionality 13
per node.
Flood occurrence labels per node were assigned as: 1 if any inventory flood point
falls within the sub-watershed, 0 otherwise.

\subsection{Baseline Machine-Learning Models}
\label{sec:baselines}

Three tree-based ensemble classifiers were trained as pixel-based baselines:
(1) Random Forest (RF; \citealt{breiman2001rf}) with 500 trees, maximum depth 6,
and balanced class weights;
(2) XGBoost \citep{chen2016xgboost} with 500 trees, maximum depth 6,
and scale\_pos\_weight\,=\,5 for class imbalance correction;
(3) LightGBM \citep{ke2017lightgbm} with 500 leaves, maximum depth 6,
and balanced class weights.
A fourth model, a stacking ensemble, used RF + XGBoost + LightGBM as base learners
and Logistic Regression as the meta-learner.
All models were implemented using scikit-learn 1.8, XGBoost 3.0, and LightGBM 4.6.

\subsection{Graph Neural Network (GraphSAGE)}
\label{sec:gnn}

The primary novel model is a three-layer GraphSAGE
\citep{hamilton2017graphsage} trained on the watershed connectivity graph.
GraphSAGE performs inductive representation learning by sampling and aggregating
feature information from a node's local neighbourhood:
\begin{equation}
  \mathbf{h}_v^{(k)} = \sigma\!\left(
    \mathbf{W}^{(k)} \cdot
    \mathrm{CONCAT}\!\left(
      \mathbf{h}_v^{(k-1)},\;
      \mathrm{AGG}\!\left(\left\{\mathbf{h}_u^{(k-1)}\,:\,u \in \mathcal{N}(v)\right\}\right)
    \right)
  \right)
  \label{eq:graphsage}
\end{equation}
where $\mathbf{h}_v^{(k)}$ is the embedding of node $v$ at layer $k$,
$\mathcal{N}(v)$ is the neighbourhood of $v$,
$\mathrm{AGG}$ is the mean aggregator,
$\mathbf{W}^{(k)}$ is a trainable weight matrix,
and $\sigma$ is a non-linearity (ReLU).

The architecture comprised three SAGEConv layers
(13 $\to$ 64 $\to$ 64 $\to$ 2 dimensions),
with ReLU activations and dropout ($p$\,=\,0.30) between layers.
Training used the Adam optimiser (learning rate $10^{-3}$, weight decay $10^{-4}$)
for 200 epochs, with a class-weighted cross-entropy loss to address
class imbalance (typically 20--30\% positive nodes).
The model was implemented in PyTorch 2.6 and PyTorch Geometric 2.7.
In settings where PyTorch Geometric is unavailable,
a neighbourhood-aggregation MLP proxy is used as a computational fallback
(Appendix~\ref{app:proxy}).

\subsection{Spatial Block Cross-Validation}
\label{sec:validation}

All models were evaluated using leave-one-basin-out (LOBO) spatial block
cross-validation with five folds corresponding to the five major river basins
(Beas, Satluj, Chenab, Ravi, Yamuna).
In each fold, the held-out basin served as the test set and the four
remaining basins as training data.
This design mirrors the intended use case: training on well-studied basins
and predicting susceptibility in a spatially separated target basin.
Reported AUC values are the mean and standard deviation across the five folds.

The following metrics were recorded per fold:
AUC-ROC (primary), F1-score (macro), Cohen's Kappa,
True Positive Rate (TPR), False Positive Rate (FPR),
and the Brier score.
All threshold-dependent metrics were computed at the default 0.5 threshold.

\subsection{Conformal Prediction}
\label{sec:conformal}

Split-conformal (inductive conformal) prediction
\citep{angelopoulos2023conformal, taquet2022mapie}
was applied to the best-performing model after fitting.
The procedure is as follows:
(1) the full training dataset is split into a proper training set (60\%),
a calibration set (20\%), and a test set (20\%);
(2) the base model is trained on the proper training set;
(3) for each calibration example $i$, the non-conformity score
$\alpha_i = 1 - \hat{p}(y_i \mid x_i)$ is computed,
where $\hat{p}(y_i \mid x_i)$ is the predicted probability of the true class;
(4) the $(1-\alpha)$-quantile of the calibration scores is used as the threshold
$\hat{q}$:
\begin{equation}
  \hat{q} = \mathrm{Quantile}\!\left(
    \{\alpha_i\}_{i=1}^{n_\mathrm{cal}},\;
    \frac{\lceil (1-\alpha)(n_\mathrm{cal}+1) \rceil}{n_\mathrm{cal}}
  \right)
  \label{eq:conformal_quantile}
\end{equation}
(5) prediction sets at coverage level $1-\alpha$ are:
$C(x) = \{y : 1 - \hat{p}(y \mid x) \leq \hat{q}\}$.

We set $\alpha = 0.10$ (90\% coverage).
Split-conformal prediction guarantees that $\mathbb{P}(y \in C(X)) \geq 1-\alpha$
for any new observation under the assumption of exchangeability
(i.e., that test examples are drawn from the same distribution as calibration examples).
Coverage was empirically verified on the held-out test set
and on the independent 2023 temporal validation set.

For spatial mapping, each pixel is assigned the point susceptibility estimate
and a lower/upper bound derived from $\hat{q}$, producing three map layers:
the point estimate, a conservative lower bound, and a conservative upper bound.
The uncertainty width $W = \hat{p}_\mathrm{upper} - \hat{p}_\mathrm{lower}$
quantifies epistemic uncertainty at each location.
MAPIE v1.3 \citep{taquet2022mapie} was used for the production implementation;
a manual fallback implementation is provided for reproducibility
(see \cref{sec:code_availability}).

\subsection{SHAP Feature Importance}
\label{sec:shap}

SHAP (SHapley Additive exPlanations; \citealt{lundberg2017shap})
TreeExplainer was applied to the RF component of the stacking ensemble
to produce both global and spatially disaggregated factor importance.
Global importance was measured as the mean absolute SHAP value per factor
across all training samples.
Spatial SHAP analysis assigned each pixel (or sub-watershed) the index of the
conditioning factor with the highest absolute SHAP value,
producing a ``dominant factor'' map that reveals geographic variation
in the primary driver of susceptibility.
District-level SHAP aggregation was computed by averaging absolute SHAP values
within each district boundary, enabling district-specific factor priority reports.
